% s-315-p315-m003; original PDF 315, 316
\Aphorism{s-315-p315-m003}{370}
Wenn es ein Fortleben nach dem Tode gäbe — welche Marter für einen ehrlichen und humanen Todten, als passiver Zuschauer seine eigene Idealisirung und Kanonisirung mit anzusehen! Hülflos Zeuge zu sein, wie die Flecken retouchirt, die Lichter verstärkt und neue aufgesetzt werden, wie die nothwendige Reibung des Zusammenlebens nachträglich als bitteres Unrecht gegen den Todten empfunden und ins Unverhältnissmässige vergrössert wird, wie aus dem Menschen der Heilige und um den Heiligen die Legende sich bildet! Ich stelle mir vor, wie Jemand, dem dies begegnete, die Gelegenheit suchen müsste, sich als \Foreign{spirit} kundzugeben, nur um mit irgend einer Rabelaisischen Derbheit zwischen diesen sentimentalen Todtencult zu fahren! Der Tod ist ja der eigentliche Beginn der Macht, die auf Verkennung beruht; der Lebende macht seinen Einfluss glücklicherweise immer wieder selbst unschädlich, er hilft den Anderen, sich über ihn zu ernüchtern. Sobald diese Fähigkeit, wider sich selbst zu wirken, aufhört, muss auch alles Andere aufhören: insbesondere die unthätige Zeugenschaft bei einer Fälschung, die, anfangs schüchtern und „unter uns“, zu einem öffentlichen Betruge — ich wollte sagen, zu einer Religion anwachsen kann. Der Todte hat ein Recht auf volles Unbewusstsein; das ist seine \Foreign{décharge} für den Fall, dass mit seinem Namen Unfug getrieben wird.
\par

